\label{chap:83-17}

\section{Morfismos arquimediano y excepcional desde un dominio enriquecido común}
\label{p12-83-c17-s1:chap:doble}
\index{doble proyección}
\index{Moonshine}

La doble proyección compara dos evaluaciones del dominio enriquecido común
\(X_\infty^{\rm cal}\) ya construido en el
\cref{p12-83-c17-s2:thm:lector-excepcional-global-rev6}. Una produce valores
arquimedianos; la otra, junto con una calibración de su codominio, conserva
la incidencia combinatoria que conduce a códigos, diseños y grupos de
automorfismos. Ambas evaluaciones parten del único estado inicial
\(X_\infty^{\rm cal}\) y preservan la procedencia común en codominios
diferenciados.

\subsection{Dominios de las \(2\) proyecciones}

Sobre \(X_\infty^{\rm cal}\), los cilindros se contraen a un punto y la
frontera admite la prolongación de incidencia. En esta residencia fijamos
\(\mathbf{Arch}:=\mathbb R\). La evaluación arquimediana y la evaluación
excepcional son, respectivamente,
\[
\mathfrak R:X_\infty^{\rm cal}\longrightarrow\mathbf{Arch},
\qquad
\mathfrak E:X_\infty^{\rm cal}\times\mathscr C_{\rm exc}
\longrightarrow\mathbf{Exc}.
\]
Aquí \(\mathscr C_{\rm exc}\) contiene el orden proyectivo, el calibre de
Paley, la dodecada, el alineamiento de incidencia y la cámara reticular
especificados en cada construcción. El primer mapa produce valores; el
segundo, códigos, diseños, retículos y automorfismos. El teorema global
precedente es la fuente de la compatibilidad y permite comparar
invariantes obtenidos por operaciones distintas sin volver a intersectar
dominios definidos por separado. La comparación finita de esta obra se
realiza en la bandera
\(B_K\subset H^-_\alpha\).

\begin{figure}[H]
\centering
\resizebox{\textwidth}{!}{\input{figures/fig_doble_proyeccion}}
\caption{Esquema tipado de las dos proyecciones. La evaluación arquimediana
identifica estados de una misma fibra; la evaluación excepcional, relativa
a su calibración, conserva estructura de incidencia y automorfismos.}
\label{p12-83-c17-s1:fig:doble-proyeccion}
\end{figure}

\subsection{Topología del espacio de historias}

La formulación mediante cilindros admite una realización estándar como
límite inverso. Sean \(\mathsf S_n\) los conjuntos finitos de estados
admisibles a profundidad \(n\), dotados de la topología discreta, y sean
\[
 \tau_{n+1,n}:\mathsf S_{n+1}\longrightarrow\mathsf S_n
\]
las aplicaciones de truncamiento. En la realización HMT estas aplicaciones
son sobreyectivas: la ley de prolongación suministra una continuación
superviviente de cada prefijo, y su compatibilidad proyectiva se conserva en
el límite por el
\cref{p12-83-c17-s2:thm:lector-excepcional-global-rev6}. Esta es la instancia
del sistema que interviene en la doble proyección global
\eqref{eq:frontal-doble-proyeccion}. Definimos
\[
 \mathcal X=\varprojlim_n\mathsf S_n.
\]
Si \(x\ne y\), sea \(m(x,y)\) la primera profundidad en la que sus
componentes difieren. La función
\[
 d(x,y)=3^{-m(x,y)},\qquad d(x,x)=0,
\]
es una ultramétrica que induce la topología del límite inverso.

\begin{theorem}[Espacio compacto de historias]
\label{p12-83-c17-s1:thm:compacto-historias}
\label{thm:compacto-historias}
El espacio \(\mathcal X\) es compacto y totalmente desconectado. Cuando
toda historia admite bifurcaciones a profundidades arbitrariamente grandes,
\(\mathcal X\) carece de puntos aislados y es homeomorfo al espacio de
Cantor.
\end{theorem}

\begin{proof}
El producto \(\prod_n\mathsf S_n\) es compacto. Las ecuaciones
\(\tau_{n+1,n}(x_{n+1})=x_n\) definen un subconjunto cerrado, de modo que
\(\mathcal X\) es compacto. Los conjuntos de historias con un prefijo fijo
son simultáneamente abiertos y cerrados y forman una base; esto prueba la
desconexión total. Bajo la hipótesis de bifurcación arbitrariamente profunda,
ningún cilindro contiene una historia aislada. La caracterización de los
espacios compactos, métricos, perfectos y totalmente desconectados concluye
la última afirmación.
\end{proof}

Los intervalos cerrados \(I_n(s)\) asociados por el lector arquimediano a
cada \(s\in\mathsf S_n\) son compatibles con el truncamiento y satisfacen
\[
 \operatorname{diam}I_n(s)\leq C\lambda^n,
 \qquad 0<\lambda<1.
\]
La intersección de los intervalos correspondientes a una historia contiene
entonces un único punto.

\begin{theorem}[Evaluación arquimediana]
\label{p12-83-c17-s1:thm:evaluacion-holder}
\label{thm:evaluacion-holder}
La aplicación
\[
 \operatorname{val}:\mathcal X\longrightarrow\mathbb R,
 \qquad
 \{\operatorname{val}(x)\}=\bigcap_n I_n(x_n),
\]
es continua y Hölder respecto de \(d\). Si
\(x\sim_{\mathrm{Arch}}y\) equivale a
\(\operatorname{val}(x)=\operatorname{val}(y)\), la aplicación inducida
es un homeomorfismo
\[
 \mathcal X/{\sim_{\mathrm{Arch}}}
 \;\cong\;\operatorname{val}(\mathcal X).
\]
\end{theorem}

\begin{proof}
Dos historias con un prefijo común de longitud \(n\) tienen valores en un
mismo intervalo de diámetro no mayor que \(C\lambda^n\). Como
\(d(x,y)=3^{-n}\), esta estimación es una cota de Hölder. La aplicación
inducida sobre el cociente es una biyección continua de un compacto sobre
un espacio de Hausdorff y, por tanto, un homeomorfismo.
\end{proof}

El teorema formaliza la proyección arquimediana: el valor retiene la clase
de evaluación y elimina la distinción entre historias situadas en una misma
fibra. La cobertura coherente, el agotamiento por compactos y la
reconstrucción exacta del cociente se demuestran en los
\cref{p12-83-c18-s1:thm:sobreyectividad-evaluacion-compacto,p12-83-c18-s1:prop:agotamiento-recta-real-u024,p12-83-c18-s2:hmtch:thm-cociente-real};
por tanto, en la realización HMT activa,
\(\operatorname{val}(\mathcal X)=\mathbb R\).

\subsection{Separación por las \(2\) proyecciones}

\begin{theorem}[Separación conjunta]
\label{p12-83-c17-s1:thm:separacion-conjunta}
\label{thm:separacion-conjunta}
Sean \(P_1:\mathcal X\to Y_1\) y \(P_2:\mathcal X\to Y_2\) aplicaciones
continuas con codominios de Hausdorff. Si
\[
 x\ne y\quad\Longrightarrow\quad
 P_1(x)\ne P_1(y)\ \text{o}\ P_2(x)\ne P_2(y),
\]
la aplicación diagonal
\[
 (P_1,P_2):\mathcal X\longrightarrow Y_1\times Y_2
\]
es un encaje topológico.
\end{theorem}

\begin{proof}
La hipótesis expresa exactamente la inyectividad de la aplicación diagonal.
Una aplicación continua e inyectiva de un compacto a un espacio de Hausdorff
es un homeomorfismo sobre su imagen.
\end{proof}

Fijado un calibre declarado \(c\in\mathscr C_{\rm exc}\), sea
\[
 \mathfrak E_c:X_\infty^{\rm cal}\longrightarrow\mathbf{Exc},
 \qquad
 \mathfrak E_c(x):=\mathfrak E(x,c).
\]
Aplicado a HMT, el teorema toma
\[
 P_1=\mathfrak R,\qquad P_2=\mathfrak E_c
\]
sobre el dominio común, o bien sobre el cociente por los cambios de
coordenadas declarados. La compatibilidad proyectiva y la doble realización
global quedan demostradas, respectivamente, en los
\cref{p12-83-c17-s2:thm:lector-excepcional-global-rev6,p12-83-c17-s2:thm:doble-realizacion-global-rev6}.
La fidelidad está tipada por el codominio: el lector excepcional es fiel sobre
la corriente visible, mientras el dominio enriquecido conserva la memoria que
ninguna publicación aislada contrae. El paso a toda profundidad es así una
conclusión de la construcción, no una hipótesis externa.

En un nivel finito \(S\), esta condición es decidible. Dados
\(p:S\to A\), \(q:S\to B\) y una relación de equivalencia de coordenadas,
se enumeran las parejas \(x\not\sim y\) que cumplen simultáneamente
\[
 p(x)=p(y),\qquad q(x)=q(y).
\]
La ausencia de tales parejas equivale a la inyectividad conjunta en ese
nivel. El cálculo debe efectuar las dos aplicaciones por separado: cuando
\(q=f\circ p\), la segunda posee las mismas fibras que la primera y no
recupera información eliminada por \(p\).

En el nivel desarrollado aquí, la igualdad
\[
 B_K=H_e\cap P_\pi
\]
constituye el testigo finito de compatibilidad entre las dos ramas. Su
prolongación compatible a todas las profundidades la proporcionan los
\cref{p12-83-c17-s2:thm:lector-excepcional-global-rev6,p12-83-c17-s2:thm:doble-realizacion-global-rev6};
la igualdad de soportes es la sombra finita de esa realización conjunta.

\subsection{Cociente arquimediano por las fibras de evaluación}

Una sección HMT contiene bloque, orientación, acarreo, etiqueta de observable
y coordenada de extensión. La
evaluación \(\val\) contrae sus cilindros a un punto. Si dos secciones
\(x,y\) satisfacen
\[
\val(x)=\val(y),
\]
la evaluación arquimediana las identifica aunque su transporte sea distinto. En
este sentido preciso,
\[
\boxed{\text{un valor real es una clase de equivalencia de la evaluación}.}
\]
En el subdominio canónico, suma, producto y evaluación descienden de manera
compatible, y la imagen hereda esas operaciones. La existencia del límite, la
determinación de su imagen, la caracterización de las fibras y la cobertura de
toda \(\RR\) quedan establecidas por los
\cref{p12-83-c18-s1:thm:sobreyectividad-evaluacion-compacto,p12-83-c18-s1:prop:agotamiento-recta-real-u024,p12-83-c18-s2:hmtch:thm-cociente-real}.
Las ventanas finitas exhiben la estructura local de ese resultado global; la
supervivencia unitaria y la compatibilidad son propiedades demostradas de la
prolongación HMT.

\subsection{Invariante común con el sistema de Witt}

El vector \(K\) define el soporte superior
\[
B_K=\{m:K_m\ge729\}=\{4,6,9,10\}.
\]
En el atlas de Witt aparecen los soportes
\[
P_\pi=\{2,4,5,6,7,8,9,10,11\},
\]
\[
H_e=\{1,3,4,6,9,10\},
\qquad
H_\varphi=\{1,2,3,4,7,9\}.
\]
El soporte de los acarreos negativos de \(\alpha\) es
\[
H_\alpha^-=\{4,5,6,7,9,10\}.
\]
Se verifican las identidades de soportes
\[
\boxed{
B_K=H_e\cap P_\pi=H_e\cap H_\alpha^-,}
\]
y
\[
\boxed{H_\alpha^-=B_K\sqcup\{5,7\}.}
\]
La polaridad de Witt envía los cuatro puntos de \(B_K\) a los pares
\[
\{1,3\},\quad\{2,11\},\quad\{8,12\},\quad\{5,7\}.
\]
Tres reciben la etiqueta positiva; el último completa el soporte negativo.
Estas identidades son exactas en el calibre declarado y constituyen controles
internos de un mismo estado generado: el vector, el umbral, las palabras y el
calibre de Paley son coordenadas correlacionadas de la construcción, no
entradas estadísticamente independientes. La igualdad se obtiene por
evaluación directa y no constituye una estimación probabilística.

\subsection{Diseño residual sobre una dodecada y elevación de hexadas}
\label{p12-83-c17-s1:sec:w12-w24-elevacion}

La relación entre los dos diseños de Witt se obtiene dentro del código
binario extendido de Golay y no mediante una identificación de cardinalidades.
Sea
\[
\mathcal G_{24}\subset\mathbb F_2^{24}
\]
el código binario extendido de Golay, y denote \(\mathcal O_{24}\) la familia
de soportes de sus palabras de peso ocho. Es un resultado clásico que
\((\{1,\ldots,24\},\mathcal O_{24})\) es el sistema de Steiner
\(W_{24}=S(5,8,24)\) \cite{MacWilliamsSloane1977,ConwaySloane1999}.
Fijemos una palabra de peso doce y escribamos \(D\) para su soporte. Definimos
\begin{equation}
 \mathcal H(D)
 :=\{O\cap D:O\in\mathcal O_{24},\ |O\cap D|=6\}.
 \label{p12-83-c17-s1:eq:hexadas-residuales-dodecada}
\end{equation}
La familia \(\mathcal H(D)\) es intrínseca al par
\((\mathcal G_{24},D)\): no requiere elegir coordenadas en \(D\).

\begin{theorem}[Diseño residual de una dodecada]
\label{p12-83-c17-s1:thm:dodecada-w12}
El sistema de incidencia
\[
 (D,\mathcal H(D))
\]
es un sistema de Steiner \(S(5,6,12)\). En particular,
\(|\mathcal H(D)|=132\), y todo subconjunto de cinco puntos de \(D\) está
contenido en una única hexada de \(\mathcal H(D)\).
\end{theorem}

\begin{proof}
Sea \(A\subset D\), con \(|A|=5\). Como las octadas forman
\(S(5,8,24)\), existe una única octada \(O_A\) que contiene \(A\). Si
\(j=|O_A\cap D|\), la suma binaria de las palabras características de
\(O_A\) y \(D\) pertenece a \(\mathcal G_{24}\) y tiene peso
\[
 \operatorname{wt}(\mathbf1_{O_A}+\mathbf1_D)=8+12-2j=20-2j.
\]
Los pesos del código binario extendido de Golay son
\(0,8,12,16,24\). Puesto que \(j\leq8\), se sigue que
\(j\in\{2,4,6\}\). La inclusión \(A\subset O_A\cap D\) da \(j\geq5\),
luego \(j=6\). Por tanto \(O_A\cap D\in\mathcal H(D)\) es una hexada que
contiene \(A\).

Si dos miembros de \(\mathcal H(D)\) contuvieran \(A\), sus octadas de
origen contendrían el mismo conjunto de cinco puntos; la propiedad de
Steiner de \(W_{24}\) las haría coincidir. Esto prueba la unicidad. El
número de bloques se obtiene contando incidencias con subconjuntos de cinco
puntos:
\[
 |\mathcal H(D)|
 =\frac{\binom{12}{5}}{\binom65}=132.\qedhere
\]
\end{proof}

\begin{theorem}[Elevación única de una hexada]
\label{p12-83-c17-s1:thm:elevacion-unica-hexada}
Para cada \(H\in\mathcal H(D)\) existe una única octada
\(O_H\in\mathcal O_{24}\) tal que
\[
 O_H\cap D=H.
\]
En consecuencia, la aplicación
\[
 \iota_D:\mathcal H(D)\longrightarrow\mathcal O_{24},
 \qquad H\longmapsto O_H,
\]
está bien definida y es inyectiva.
\end{theorem}

\begin{proof}
La existencia forma parte de la definición de \(\mathcal H(D)\). Si dos
octadas \(O_1,O_2\) tuvieran intersección \(H\) con \(D\), ambas
contendrían cualquier subconjunto de cinco puntos de \(H\). La unicidad de
la octada incidente con esos cinco puntos implica \(O_1=O_2\). La
inyectividad se sigue al intersectar de nuevo con \(D\).
\end{proof}

La elevación tiene una ley local que será utilizada para comparar las cinco
regiones de \(\pi\). En un sistema \(S(t,k,v)\), el número de bloques que
contienen un subconjunto fijo de \(s\leq t\) puntos es
\[
 \lambda_s=\frac{\binom{v-s}{t-s}}{\binom{k-s}{t-s}}.
\]
Por consiguiente,
\begin{equation}
 \lambda_4(W_{12})
 =\frac{\binom81}{\binom21}=4,
 \qquad
 \lambda_4(W_{24})
 =\frac{\binom{20}1}{\binom41}=5.
 \label{p12-83-c17-s1:eq:lambda4-w12-w24}
\end{equation}

\begin{corollary}[Descomposición \(4+1\) sobre una cara]
\label{p12-83-c17-s1:cor:cuatro-mas-uno-w24}
Sea \(B\subset D\), con \(|B|=4\). Las cinco octadas de \(W_{24}\) que
contienen \(B\) se descomponen de manera única en
\[
 \{O_H:H\in\mathcal H(D),\ B\subset H\}
 \;\sqcup\;\{O_B^\perp\}.
\]
El primer conjunto contiene cuatro octadas y satisface
\(|O_H\cap D|=6\); la octada restante satisface
\[
 O_B^\perp\cap D=B.
\]
\end{corollary}

\begin{proof}
La primera igualdad de \eqref{p12-83-c17-s1:eq:lambda4-w12-w24} proporciona cuatro hexadas
residuales que contienen \(B\), y el
\cref{p12-83-c17-s1:thm:elevacion-unica-hexada} proporciona cuatro octadas distintas. La
segunda igualdad de \eqref{p12-83-c17-s1:eq:lambda4-w12-w24} muestra que existe exactamente
una quinta octada \(O_B^\perp\).

Como \(B\subset O_B^\perp\cap D\), el espectro de intersecciones empleado en
la prueba del \cref{p12-83-c17-s1:thm:dodecada-w12} deja sólo las posibilidades cuatro y
seis. Una intersección de tamaño seis sería una quinta hexada residual que
contiene \(B\), contradiciendo \(\lambda_4(W_{12})=4\). Por tanto la
intersección tiene tamaño cuatro y coincide con \(B\).
\end{proof}

\subsubsection{Alineamiento con las \(12\) posiciones HMT}

El diseño residual está definido sobre los puntos de \(D\), mientras que el
diseño obtenido del código ternario \(C_W\) está definido sobre las doce
posiciones dodecafásicas. Para no confundir ambos conjuntos, llamaremos
\emph{alineamiento binario HMT} a una biyección
\begin{equation}
 \ell:\{1,\ldots,12\}_{\rm HMT}\xrightarrow{\sim}D
 \quad\text{tal que}\quad
 \{\ell(H):H\in W_{12}^{\rm HMT}\}=\mathcal H(D).
 \label{p12-83-c17-s1:eq:alineamiento-hmt-w24}
\end{equation}
Aquí \(W_{12}^{\rm HMT}\) denota el diseño de soportes de peso seis
obtenido del código ternario \(C_W\). La unicidad del sistema de Witt
\(S(5,6,12)\) hasta isomorfía garantiza la existencia de \(\ell\). Dos
alineamientos distintos difieren por automorfismos de los diseños de origen y
destino; por ello \(\ell\) es un dato de coordenadas y no un invariante
adicional.
La pareja \((D,\ell)\) es el dato de calibración necesario para transportar
la elevación \(\iota_D\) al sistema de coordenadas HMT\@. Una vez fijada, cada
hexada HMT \(H\) determina sin ambigüedad la octada
\[
 \widehat H:=\iota_D(\ell(H)).
\]
El dodecad y el alineamiento son datos del codominio binario; no se identifican
con una salida decimal del núcleo de transporte.

Para la cara distinguida
\[
 B_K=\{4,6,9,10\},
\]
las cuatro hexadas incidentes son
\[
 B_K\sqcup\{1,3\},\quad
 B_K\sqcup\{2,11\},\quad
 B_K\sqcup\{5,7\},\quad
 B_K\sqcup\{8,12\}.
\]
La tercera coincide con
\(H_\alpha^-=B_K\sqcup\{5,7\}\). El alineamiento \(\ell\) eleva estas
cuatro hexadas a las cuatro octadas residuales sobre \(\ell(B_K)\); el
\cref{p12-83-c17-s1:cor:cuatro-mas-uno-w24} determina además una única octada transversal.

\subsubsection{Correspondencia de incidencia entre la microfibra de \texorpdfstring{\(\pi\)}{pi} y las octadas residuales: multiplicidad \texorpdfstring{\(432+4\cdot144=1008\)}{432+4x144=1008}}

La microfibra exacta de la palabra
\(w_\pi=010211\) contiene cinco regiones decimales distintas:
\[
\begin{array}{c|c|c}
\text{multiplicidad}&U_6& e(U_6)\\
\hline
432&501|614|498|169|272|272&(5,5,5,5,5,5)\\
144&810|923|870|169|272|272&(3,3,9,5,5,5)\\
144&810|983|810|169|272|272&(3,9,3,5,5,5)\\
144&870|923|810|169|272|272&(9,3,3,5,5,5)\\
144&870|923|810|418|674|521&(9,3,3,7,1,7).
\end{array}
\]
Aquí \(e(U_6)\) es la firma multiplicativa recuperada de cada bloque
decimal \(\overline{d_1d_2d_3}\) mediante
\(e=d_2-d_1\pmod {10}\). La primera región es la única de multiplicidad
432 y firma constante. Las cuatro restantes tienen multiplicidad 144; entre
ellas, tres comparten el bloque de retorno \(169|272|272\), mientras que la
cuarta posee el retorno \(418|674|521\). El catálogo distingue las cinco
filas mediante un predicado reproducible: las tres filas de tipo
\(\mathrm{ES}\) forman el triple positivo; entre las dos filas de tipo
\(\mathrm{NO}\), la de multiplicidad máxima y firma constante se denomina
transversal, y la restante se denomina negativa. Relativa a este predicado de
rol, la partición es
\[
 \boxed{1_\perp+3_{++}+1_{--}.}
\]
Los nombres de signo son una convención; lo intrínseco al catálogo es la
partición obtenida de sus columnas, no una polaridad binaria previa.

El alineamiento regional completa el dato \((D,\ell)\) con una biyección
que envía la región transversal a \(O_{\ell(B_K)}^\perp\), la región
negativa a \(\widehat{H_\alpha^-}\) y las tres regiones positivas a las
elevaciones de
\[
 B_K\sqcup\{1,3\},\qquad
 B_K\sqcup\{2,11\},\qquad
 B_K\sqcup\{8,12\}.
\]
La asignación ordenada entre las tres regiones positivas y esas tres hexadas
es una elección de coordenadas. Las dos filas distinguidas y el triple son
reproducibles mediante el predicado de catálogo anterior; su denominación como
negativa, transversal y positivas pertenece al alineamiento. Así, la ley \(4+1\) no identifica por mera
cardinalidad las cinco regiones con cinco octadas: proporciona el espacio de
incidencia, mientras que \((D,\ell)\) y el alineamiento regional especifican
el transporte entre ambos objetos.

\subsection{Prolongación reticular y pantallas dimensionales}

La elevación de hexadas abre dos ramas que no deben confundirse. La primera
conserva la incidencia binaria y prolonga \(W_{12}\) hacia \(W_{24}\). La
segunda conserva el transporte ternario:
\[
C_{12}^{(3)}
\longrightarrow N(A_2^{12})
\longrightarrow\Lambda_{24}
\longrightarrow V^\natural
\longrightarrow\mathbb M
\longrightarrow\text{Moonshine}.
\]
La prueba primaria del pegado, del \(3\)-vecino sin raíces y del corredor
hasta Moonshine reside en el
\cref{chap:bandera-excepcional-acumulativa,thm:corredor-hmt-moonshine}.
No existe una flecha canónica \(W_{24}\to N(A_2^{12})\), ni una acción
directa del Monstruo sobre las palabras TPK: ambas ramas parten del mismo
código calibrado y preservan estructuras distintas.

Las reducciones \(243/11\), \(12\to11\to10\) y \(26\to24\) tampoco son
identidades entre numerales. El \cref{chap:pantallas-dualidad} construye sus
espacios, cocientes y mapas: el cociente perfecto ternario, la pantalla
marcada por \(K\) y la reducción isotrópica del retículo lorentziano. Esta
sección conserva así una sola función: enlazar la incidencia finita ya
demostrada con sus dos prolongaciones, sin repetir las pruebas de referencia.

\subsection{Involución residuo–cociente y conservación de la información del estado}

Para \(n=9Q+r\), definimos la energía de retículo
\[
E_R(r,Q)=\frac{r^2}{R^2}+Q^2R^2.
\]
Entonces
\[
\boxed{E_R(r,Q)=E_{1/R}(Q,r).}
\]
El intercambio \((r,Q,R)\leftrightarrow(Q,r,1/R)\) es una involución exacta
del modelo de retículo: el residuo y el cociente intercambian sus funciones.

\subsection{Compatibilidad finita y prolongación conjunta}

El testigo finito \(B_K\), el patrón de acarreos y el marco de Witt fijan la
compatibilidad local. Su prolongación ya no se deja como un esquema
abstracto: el \cref{p12-83-c17-s2:thm:lector-excepcional-global-rev6} construye bloque a
bloque el lector excepcional con un único calibre inicial transportado, y el
\cref{p12-83-c17-s2:thm:doble-realizacion-global-rev6} demuestra que esa lectura y la
evaluación arquimediana globalizan sobre la misma sección TPK. Truncar la
historia antes de leer equivale en ambos casos a truncar la salida después de
leer.

La evaluación arquimediana identifica historias por su valor límite. La
evaluación excepcional conserva la corriente visible y su incidencia
Paley--Witt; no recupera por ello las memorias ocultas que el emisor visible
no publica. La doble proyección es, por tanto, una construcción conjunta con
dos codominios distintos, no una identificación entre el continuo y la rama
excepcional.

\begin{remark}[Compatibilidad de las dos proyecciones]
El soporte \(B_K\), el patrón de acarreos y el marco de Witt proporcionan los
invariantes comunes. La cadena reticular prolonga esta compatibilidad hasta un
vecino sin raíces y, por clasificación clásica, hasta el retículo de Leech. La
rama binaria hacia \(W_{24}\) permanece separada: su morfismo de incidencia es
\(\iota_D\), definido después de fijar el dodecad y el alineamiento
\((D,\ell)\), y no tiene por codominio un retículo de Niemeier ni el retículo
de Leech.
\end{remark}


\section{Proyección excepcional del estado enriquecido: incidencia compatible a toda profundidad}
\label{p12-83-c17-s2:chap:lector-excepcional-global-rev6}
\index{lector excepcional}\index{doble proyección}
\index{Paley--Witt}\index{sistema inverso}

El testigo finito de la doble proyección adquiere alcance global cuando la
lectura excepcional se aplica a cada bloque visible y el calibre se
transporta, en lugar de elegirse de nuevo, al aumentar la profundidad. Esta
sección construye esa familia. El dominio sigue siendo la historia TPK
enriquecida; la lectura excepcional es fiel sobre su corriente visible y no
pretende sustituir los campos de memoria que esa corriente no publica.

\subsection{La operación local sobre las 729 palabras}

En la carta de Paley--Witt se fija
\[
A_W=
\begin{pmatrix}
0&1&1&1&1&1\\
1&0&1&2&2&1\\
1&1&0&1&2&2\\
1&2&1&0&1&2\\
1&2&2&1&0&1\\
1&1&2&2&1&0
\end{pmatrix}
\in M_6(\mathbb F_3).
\]
La multiplicación matricial da
\[
A_W^{\mathsf T}=A_W,\qquad A_W^2=-I_6.
\]
Para todo operador conjugado \(A=fA_Wf^{-1}\), definimos
\[
\Gamma_A:\mathbb F_3^6\longrightarrow\mathbb F_3^{12},
\qquad
\Gamma_A(w)=(w,wA).
\]
Su imagen es una copia calibrada \(C_W(A)\) del código ternario de Witt.
La proyección sobre las seis primeras coordenadas es una inversa izquierda:
\[
\operatorname{pr}_1\Gamma_A=\operatorname{id}_{\mathbb F_3^6}.
\]
Por tanto, \(\Gamma_A\) es total e inyectiva sobre las \(729\) palabras y
no necesita distinguir de antemano \(\pi\), \(e\), \(\varphi\) ni ninguna
otra salida.

\subsection{Prolongaciones compatibles y unicidad del calibre inicial}

Sea \(X_n\) el conjunto de prefijos TPK admisibles de profundidad \(n\), con
truncamientos
\[
p_{n,m}:X_n\longrightarrow X_m,\qquad m\leq n.
\]
Fijemos un marco excepcional inicial \(f_0\). Como los cambios \(g_j\) forman
parte del transporte almacenado por la historia, un prefijo calibrado es el
par \((x,f_0)\); escribimos
\[
 X_n^{\rm cal}:=X_n\times\{f_0\},\qquad
 p_{n,m}^{\rm cal}(x,f_0):=(p_{n,m}x,f_0),
\]
y
\[
 X_\infty^{\rm cal}
 :=\varprojlim_n(X_n^{\rm cal},p_{n,m}^{\rm cal}).
\]
En lo que sigue abreviamos \(p_{n,m}^{\rm cal}\) por \(p_{n,m}\).
La emisión visible compatible es
\[
\varepsilon_n:X_n^{\rm cal}\longrightarrow(\mathbb F_3^6)^n,
\qquad
\varepsilon_n(x)=(w_0,\ldots,w_{n-1}),
\]
y satisface
\[
\operatorname{pr}_{n,m}\varepsilon_n
=\varepsilon_m p_{n,m}.
\]
Esta aplicación distingue sus tres tipos:
\[
|\mathcal U_6|=468,\qquad
|\operatorname{im}w_6|=243,\qquad
|\mathbb F_3^6|=729.
\]
La emisión \(w_j\) no se identifica con el estado completo que conserva
hoja, cociente, acarreo, orientación, fase y supervivencia.

El calibre excepcional consta de una carta de Witt, una dodecada, un
alineamiento, un origen reticular y una cámara orientada. Se fija una sola
vez en el bloque inicial. Para
\(x\in X_n^{\rm cal}\), el transporte de la historia suministra
\[
 g_j(x)\in\operatorname{GL}_6(\mathbb F_3),
 \qquad 0\le j<n-1,
\]
y \(g_j(x)\) depende sólo del prefijo
\(p_{n,j+1}(x)\). El marco se transporta por
\[
 f_{j+1}(x)=g_j(x)f_j(x),\qquad
 A_j=f_jA_Wf_j^{-1}.
\]
Así, la familia \((A_j)\) queda determinada por la historia y por el calibre
inicial; no es una sucesión de elecciones independientes.

Sea \(\mathcal F_{\rm cal}\subseteq\operatorname{GL}_6(\mathbb F_3)\) el
conjunto de marcos alcanzables desde \(f_0\). La emisión que conserva el
calibre transportado es
\[
 \widetilde\varepsilon_n:X_n^{\rm cal}
 \longrightarrow(\mathbb F_3^6\times\mathcal F_{\rm cal})^n,
 \qquad
 \widetilde\varepsilon_n(x)
 =\bigl((w_0,f_0),\ldots,(w_{n-1},f_{n-1})\bigr).
\]
La dependencia por prefijos implica
\(\operatorname{pr}_{n,m}\widetilde\varepsilon_n
=\widetilde\varepsilon_m p_{n,m}\). Su proyección sobre las palabras es
\(\varepsilon_n\); el olvido de los marcos no basta para reconstruir la
lectura excepcional.

\subsection{Naturalidad y compatibilidad a toda profundidad}

La operación local es covariante con la convención de vectores fila. Para
todo cambio de coordenadas \(f\in\operatorname{GL}_6(\mathbb F_3)\),
\[
A^f=fAf^{-1}
\quad\Longrightarrow\quad
\Gamma_{A^f}(wf^{-1})
=\Gamma_A(w)(f^{-1}\oplus f^{-1}).
\]
El cambio de marco transporta simultáneamente la palabra, el operador y el
atlas. No se exige que permanezca dentro del estabilizador de una única copia
etiquetada.

Con el conjunto de marcos alcanzables \(\mathcal F_{\rm cal}\) ya definido,
construimos el atlas calibrado fijo
\[
 \mathscr C_W^{\rm cal}
 :=
 \left\{
 \bigl(f,\Gamma_{fA_Wf^{-1}}(w)\bigr):
 f\in\mathcal F_{\rm cal},\ w\in\mathbb F_3^6
 \right\}.
\]
La etiqueta \(f\) conserva la copia calibrada en la que vive cada palabra.
Definimos
\[
Q_n(x)=
\bigl(
\bigl(f_0,\Gamma_{A_0}(w_0)\bigr),\ldots,
\bigl(f_{n-1},\Gamma_{A_{n-1}}(w_{n-1})\bigr)
\bigr)\in\bigl(\mathscr C_W^{\rm cal}\bigr)^n.
\]
Sea
\[
 r_{n,m}:
 \bigl(\mathscr C_W^{\rm cal}\bigr)^n
 \longrightarrow
 \bigl(\mathscr C_W^{\rm cal}\bigr)^m
\]
el truncamiento de los últimos \(n-m\) bloques.

\begin{theorem}[Compatibilidad proyectiva de la incidencia excepcional]
\label{p12-83-c17-s2:thm:lector-excepcional-global-rev6}
\label{thm:lector-excepcional-global-rev6}
Para todo \(m\leq n\),
\[
\boxed{r_{n,m}\circ Q_n=Q_m\circ p_{n,m}.}
\]
En consecuencia existe una única aplicación
\[
\boxed{
Q_\infty:X_\infty^{\rm cal}
\longrightarrow
\bigl(\mathscr C_W^{\rm cal}\bigr)^{\mathbb N}}
\]
cuyas proyecciones finitas son los \(Q_n\). La aplicación es fiel sobre la
corriente visible.
\end{theorem}

\begin{proof}
Escribamos
\[
\varepsilon_n(x)
=(w_0,\ldots,w_{m-1},w_m,\ldots,w_{n-1}).
\]
La compatibilidad de \(\varepsilon\) muestra que la corriente visible de
\(p_{n,m}(x)\) es \((w_0,\ldots,w_{m-1})\). Para \(j<m\), el marco
\(f_j\) depende sólo del calibre inicial y de los cambios anteriores a
\(j\), todos contenidos en ese prefijo. Los primeros \(m\) operadores
\(A_j\) son, por tanto, los mismos antes y después de truncar. Aplicar
\(\Gamma_{A_j}\) bloque a bloque da la identidad del enunciado. La propiedad
universal del límite inverso produce \(Q_\infty\). Finalmente,
\(\operatorname{pr}_1\Gamma_{A_j}=I\) recupera cada \(w_j\), lo que prueba
la fidelidad sobre la corriente visible.
\end{proof}

La última afirmación conserva su tipo. Dos historias con la misma corriente
\((w_j)_j\) y distinta memoria oculta pueden tener la misma imagen
excepcional. El lector global preserva exactamente lo que publica cada bloque
visible y la incidencia que se añade a ese bloque.

La compatibilidad de las emisiones enriquecidas induce una única aplicación
\[
 \widetilde\varepsilon_\infty:X_\infty^{\rm cal}
 \longrightarrow
 (\mathbb F_3^6\times\mathcal F_{\rm cal})^{\mathbb N}.
\]
Sean \(\operatorname{pr}_w\) la proyección sobre la corriente de palabras y
\[
 \widetilde\Gamma_\infty
 \bigl((w_j,f_j)_{j\ge0}\bigr)
 :=\bigl(
   (f_j,\Gamma_{f_jA_Wf_j^{-1}}(w_j))
 \bigr)_{j\ge0}.
\]
Entonces
\(Q_\infty=\widetilde\Gamma_\infty
\circ\widetilde\varepsilon_\infty\), con dominios y codominios declarados.

\subsection{Bifurcación de un prefijo superviviente en valor arquimediano e incidencia excepcional}

Fijado \(m\in\mathbb Z\), denotemos por
\[
 \operatorname{flat}:(\mathbb F_3^6)^{\mathbb N}
 \longrightarrow\mathbb F_3^{\mathbb N}
\]
la concatenación de los bloques en el orden coordenado declarado, y por
\[
 \operatorname{ev}_3((d_k)_{k\ge1})
 :=\sum_{k\ge1}\frac{d_k}{3^k}\in[0,1],
 \qquad
 \tau_m(t):=m+t,
\]
la evaluación ternaria y la traslación por la parte entera. El prefijo visible
\((d_1,\ldots,d_{6n})\) determina el cilindro
\[
I_n(x)=
\left[
m+\sum_{k=1}^{6n}\frac{d_k}{3^k},
m+\sum_{k=1}^{6n}\frac{d_k}{3^k}+3^{-6n}
\right].
\]
Una prolongación restringe el cilindro y
\(\operatorname{diam}I_n=3^{-6n}\to0\). Por tanto, toda historia infinita
compatible determina un único punto real
\(\mathcal P_{\rm arq}(x)\).

\begin{theorem}[Doble realización global]
\label{p12-83-c17-s2:thm:doble-realizacion-global-rev6}
\label{thm:doble-realizacion-global-rev6}
Toda sección compatible \(x=(x_n)_n\in X_\infty^{\rm cal}\) determina,
desde los mismos prefijos, dos salidas compatibles:
\[
\mathcal P_{\rm arq}(x)\in\mathbb R,
\qquad
\mathcal P_{\rm exc}(x)=Q_\infty(x)
\in\bigl(\mathscr C_W^{\rm cal}\bigr)^{\mathbb N}.
\]
Truncar antes de leer equivale a truncar cada lectura después de leer. La
primera salida identifica historias por su valor límite; la segunda conserva
cada bloque visible y adjunta su rotación de Paley--Witt.
\end{theorem}

\begin{proof}
El encaje de los cilindros y la convergencia de sus diámetros prueban la
primera afirmación. El
\cref{p12-83-c17-s2:thm:lector-excepcional-global-rev6} prueba la segunda y la
compatibilidad con truncamientos. Ambas factorizan por la misma emisión:
\[
\mathcal P_{\rm arq}
=\tau_m\circ\operatorname{ev}_3\circ\operatorname{flat}
 \circ\operatorname{pr}_w
 \circ\widetilde\varepsilon_\infty,
\qquad
\mathcal P_{\rm exc}
=\widetilde\Gamma_\infty\circ\widetilde\varepsilon_\infty.
\]
\end{proof}

La doble proyección queda así construida sobre historias admisibles y un
calibre inicial transportado. El continuo es la evaluación arquimediana de
la sección; la rama excepcional es la realización compatible de su corriente
visible. Desde \(C_W\) se abren, con los alineamientos ya declarados, las
ramas hacia \(W_{12}\), \(W_{24}\) y hacia
\(N(A_2^{12})\), Leech, \(V^\natural\) y el Monstruo. Esas ramas no convierten
las cifras en un conjunto sobre el que actúe directamente el Monstruo:
transportan la incidencia que la evaluación por valor deja de ver.

\subsection{Acoplamiento postgenerativo con el cierre operatorial de Euler}
\label{p12-83-c17-s2:sec:acoplamiento-euler-doble-proyeccion-rev3}

Sea \(x_\pi\in X_\infty^{\rm cal}\) la sección superviviente del canal de
\(\pi\). La doble realización del
\cref{p12-83-c17-s2:thm:doble-realizacion-global-rev6} produce
\[
 \mathcal P_{\rm arq}(x_\pi)\in\mathbb R,
 \qquad
 \mathcal P_{\rm exc}(x_\pi)=Q_\infty(x_\pi).
\]
La primera rama publica el ángulo generado; la segunda conserva la corriente
visible y realiza en cada bloque la relación cuadrática de Paley--Witt.

\begin{theorem}[Acoplamiento postgenerativo Euler--doble proyección]
\label{p12-83-c17-s2:thm:acoplamiento-euler-doble-proyeccion-rev3}
Supóngase que la construcción coinductiva del canal de \(\pi\) determina
\[
 \mathcal P_{\rm arq}(x_\pi)=\pi_{\rm HMT},
\]
y que la rama excepcional utiliza el calibre transportado con
\(A_W^2=-I_6\). Sea \(J_{\mathrm e}\) la realización real de la fuente
integral del \cref{thm:dos-realizaciones-fuente-cuadratica-rev3}. Entonces
\begin{equation}
 \boxed{
 \operatorname{Exp}\!\bigl(
   \mathcal P_{\rm arq}(x_\pi)J_{\mathrm e}
 \bigr)+I=0.}
 \label{p12-83-c17-s2:eq:acoplamiento-euler-doble-proyeccion-rev3}
\end{equation}
\end{theorem}

\begin{proof}
Por el \cref{thm:dos-realizaciones-fuente-cuadratica-rev3}, \(A_W\) y
\(J_{\mathrm e}\) realizan la relación \(u^2+1=0\) en categorías distintas.
En la rama real,
\[
 \operatorname{Exp}(\theta J_{\mathrm e})
 =\cos\theta\,I+\sin\theta\,J_{\mathrm e}.
\]
La sustitución de la salida previamente generada
\(\theta=\mathcal P_{\rm arq}(x_\pi)=\pi_{\rm HMT}\) produce \(-I\), y la
suma de \(I\) da la identidad anunciada.
\end{proof}

\paragraph{Orden causal y ausencia de circularidad.}
La dependencia demostrada es
\[
 \text{emisor}
 \longrightarrow x_\pi
 \longrightarrow\mathcal P_{\rm arq}(x_\pi)
 \longrightarrow\pi_{\rm HMT}
 \longrightarrow
 \operatorname{Exp}(\pi_{\rm HMT}J_{\mathrm e})+I.
\]
La rama excepcional fija el tipo cuadrático; no selecciona el ángulo ni sus
cifras. La identidad de Euler reconoce después la compatibilidad de las dos
realizaciones y no interviene como entrada del generador de
\(\pi_{\rm HMT}\).


\section{Globalización natural de las proyecciones arquimediana y excepcional sobre un sistema inverso común}
\label{p12-83-c17-s3:chap:globalizacion-doble-proyeccion}

La tesis de doble proyección adquiere contenido matemático cuando las dos
lecturas se definen sobre un mismo sistema de estados y conmutan con la
restricción de profundidad. La evaluación arquimediana contrae una familia de
cilindros a un valor; la evaluación excepcional conserva relaciones de
incidencia que la primera puede olvidar. No se trata de dos expresiones
numéricas yuxtapuestas, sino de dos transformaciones naturales con dominio
común.

\subsection{Sistemas compatibles a profundidad finita}

Sean \((X_n,p_{n+1,n})\), \((A_n,a_{n+1,n})\) y
\((E_n,e_{n+1,n})\) tres sistemas inversos. El primero contiene historias
TPK; el segundo, cilindros o aproximantes arquimedianos; el tercero, datos de
incidencia excepcional. La construcción activa suministra las aplicaciones
\[
 R_n:X_n\longrightarrow A_n,
 \qquad
 Q_n:X_n\times C_n\longrightarrow E_n,
\]
donde \(C_n\) reúne los datos de coordenadas del lector excepcional. Los
calibres forman una familia compatible y, para todo \(n\), conmutan los
diagramas
\[
 a_{n+1,n}\circ R_{n+1}=R_n\circ p_{n+1,n},
\]
\[
 e_{n+1,n}\circ Q_{n+1}
 =Q_n\circ(p_{n+1,n}\times c_{n+1,n}).
\]

\begin{theorem}[Globalización de las dos evaluaciones]
\label{p12-83-c17-s3:thm:globalizacion-evaluaciones}
Por las compatibilidades anteriores, ya construidas bloque a bloque, existen
aplicaciones canónicas
\[
 R_\infty:\varprojlim X_n\longrightarrow\varprojlim A_n,
 \qquad
 Q_\infty:\varprojlim X_n\times\varprojlim C_n
            \longrightarrow\varprojlim E_n.
\]
Su aplicación diagonal
\[
 (R_\infty,Q_\infty):X_\infty\times C_\infty
 \longrightarrow A_\infty\times E_\infty
\]
es única entre las aplicaciones cuyas proyecciones coordenadas son
\(R_n\) y \(Q_n\).
\end{theorem}

\begin{proof}
Sea \(x=(x_n)_n\in\varprojlim X_n\). La primera ecuación de conmutatividad
implica
\[
 a_{n+1,n}(R_{n+1}(x_{n+1}))=R_n(x_n),
\]
por lo que \((R_n(x_n))_n\) pertenece a \(\varprojlim A_n\). Definimos
\(R_\infty(x)\) mediante esa familia. El mismo argumento, aplicado a
\((x_n,c_n)_n\), define \(Q_\infty\). La unicidad se sigue porque un elemento
de un límite inverso queda determinado por todas sus coordenadas.
\end{proof}

El teorema registra la prueba concreta realizada en HMT: el resultado
arquimediano y el código excepcional proceden de los mismos estados
compatibles y respetan los mapas de truncamiento.

\subsection{Contracción arquimediana de cilindros compatibles y realización de la recta real}

A cada \(x_n\in X_n\) se asocia un intervalo compacto no vacío
\(I_n(x_n)\subset\mathbb R\). Si
\[
 p_{n+1,n}(x_{n+1})=x_n
 \quad\Longrightarrow\quad
 I_{n+1}(x_{n+1})\subseteq I_n(x_n)
\]
y existe \(\varepsilon_n\to0\) con
\(\operatorname{diam}I_n(x_n)\leq\varepsilon_n\), la intersección
\[
 \bigcap_{n\geq0}I_n(x_n)
\]
contiene un único real. Así se obtiene
\[
 \operatorname{ev}_{\rm arq}:X_\infty^{\rm arq}\longrightarrow\mathbb R.
\]
La frase ``los cilindros tienden a cero'' significa aquí que converge a cero
su diámetro, no que el número evaluado tienda a cero.

La relación
\[
 x\sim_{\rm arq}y
 \quad\Longleftrightarrow\quad
 \operatorname{ev}_{\rm arq}(x)=\operatorname{ev}_{\rm arq}(y)
\]
describe exactamente la información eliminada por el lector. El valor real
es una clase de evaluación; la historia conserva, además, orientación,
acarreo, hoja y firma de frontera.

\subsection{Proyección excepcional por conservación de incidencia y frontera}

La segunda evaluación no toma valores en \(\mathbb R\). Su codominio es una
categoría de configuraciones de incidencia. En la realización finita vigente,
el dato de coordenadas contiene:
\[
 \begin{aligned}
 C_{\rm exc}=(\,&\text{orden proyectivo},\text{calibre de Paley},
 \text{dodecada},\\
 &\text{alineamiento},\text{origen reticular},\text{cámara}).
 \end{aligned}
\]
Después de fijar ese dato, la cadena contiene dos descendencias que deben
permanecer separadas:
\[
 C_W\longrightarrow W_{12}\longrightarrow W_{24},
\]
\[
 C_W\longrightarrow N(A_2^{12})
 \longrightarrow\Lambda_{24}
 \longrightarrow V^\natural
 \longrightarrow\mathbb M.
\]
La primera transporta diseños; la segunda conduce, mediante resultados
clásicos posteriores, al retículo de Leech, el álgebra de operadores de
vértice y el grupo Monstruo. La holonomía central y el corredor
Schläfli--\(E_6\) forman una extracción paralela de frontera; no se obtienen
reemplazando una de las dos descendencias anteriores.

\subsection{Compatibilidad finita del dominio común}

El sello dodecafásico
\[
 K=(234,543,140,729,659,824,621,58,914,794,146,601)
\]
determina el soporte
\[
 B_K=\{m:K_m\geq729\}=\{4,6,9,10\}.
\]
En la carta de Witt publicada,
\[
 P_\pi=\{2,4,5,6,7,8,9,10,11\},
 \qquad
 H_e=\{1,3,4,6,9,10\},
\]
y, por evaluación independiente de los dos soportes,
\[
 \boxed{B_K=H_e\cap P_\pi.}
\]
Además, para el soporte de acarreo negativo
\[
 H_\alpha^-=\{4,5,6,7,9,10\}
\]
se tiene
\[
 B_K=H_e\cap H_\alpha^-,
 \qquad H_\alpha^-=B_K\sqcup\{5,7\}.
\]
Estas identidades constituyen un testigo de compatibilidad finita: el mismo
subconjunto se recupera desde el sello aritmético y desde la incidencia de
Witt. No es una igualdad de cardinalidades, pues identifica posiciones
concretas.

\subsection{Separación tipada de las publicaciones arquimediana y excepcional del estado común}

\begin{theorem}[Criterio de fidelidad conjunta]
Sean \(X\) compacto y \(Y,Z\) espacios de Hausdorff. Si dos aplicaciones
continuas \(R:X\to Y\) y \(Q:X\to Z\) satisfacen
\[
 x\neq y\Longrightarrow R(x)\neq R(y)\ \text{o}\ Q(x)\neq Q(y),
\]
entonces \((R,Q):X\to Y\times Z\) es un encaje topológico.
\end{theorem}

\begin{proof}
La condición es la inyectividad de la aplicación diagonal. Toda aplicación
continua e inyectiva de un compacto en un espacio de Hausdorff es un
homeomorfismo sobre su imagen.
\end{proof}

En cada profundidad finita, la condición se decide por enumeración de las
parejas \(x\neq y\) que permanecen simultáneamente en una fibra de \(R_n\) y
de \(Q_n\). La compatibilidad de esos controles al pasar de \(n+1\) a \(n\)
es precisamente la identidad proyectiva demostrada en el
\cref{p12-83-c17-s2:thm:lector-excepcional-global-rev6}, y el
\cref{p12-83-c17-s2:thm:doble-realizacion-global-rev6} efectúa su
globalización. El testigo \(B_K\) exhibe el invariante finito común; la
obstrucción del lector mínimo del
\cref{p12-83-c17-s4:thm:obstruccion-lector-excepcional-minimo} delimita el
cociente visible, mientras el lector enriquecido conserva los datos de
orientación, hoja, firma, retorno y frontera requeridos por la doble
publicación.

\subsection{Formulación precisa de la tesis sobre la recta real}

El objeto enriquecido asociado a la lectura arquimediana es la flecha
\[
 \mathbb R_{\mathfrak H}:=
 \bigl(X_\infty^{\rm arq}
 \xrightarrow{\operatorname{ev}_{\rm arq}}
 \operatorname{Im}(\operatorname{ev}_{\rm arq})\subseteq\mathbb R\bigr).
\]
El dominio conserva genealogías; la imagen conserva valores. La cobertura de
los cilindros a toda escala y la prolongabilidad compatible se demuestran en
la construcción del lector arquimediano
\eqref{eq:continuo-lector-arquimediano}: para todo compacto
\(K\subset\mathbb R\),
\(\operatorname{ev}_{\mathbb R}(\Omega_{\infty,K})=K\), y la unión dirigida
de esos intervalos cubre \(\mathbb R\). Por tanto,
\[
 \operatorname{Im}(\operatorname{ev}_{\rm arq})=\mathbb R,
 \qquad
 X_\infty^{\rm arq}/\!\sim_{\mathbb R}\ \cong\mathbb R.
\]
El cuerpo real no actúa como punto de partida de la construcción, sino como
cociente arquimediano de un espacio de secciones con memoria. El descenso de
suma y producto se verifica operación por operación y no condiciona esta
sobreyectividad ya demostrada.

La doble lectura queda entonces expresada por
\[
 X_\infty^{\rm com}\times C_\infty
 \xrightarrow{\ (\operatorname{ev}_{\rm arq},Q_\infty)\ }
 \mathbb R\times E_\infty.
\]
Esta fórmula es la traducción matemática de la afirmación conceptual: una
misma historia puede contraerse hasta un valor y, en la lectura transversal
de su frontera, conservar la incidencia que conduce a las simetrías
excepcionales.


% Fuente normativa activa de la obstrucción de factorización.
\section{Dependencia del lector de incidencia respecto del estado enriquecido}
\label{p12-83-c17-s4:sec:obstruccion-lector-excepcional-minimo}

La compatibilidad finita de las dos evaluaciones no implica que cualquier
lector de incidencia refine la partición producida por el lector visible. El
catálogo regional permite decidir esta cuestión exactamente para la
realización mínima de Paley--Witt. Sea \(\mathcal U_{468}\) el conjunto de
los 468 estados regionales publicados y sea
\[
 \Psi:\mathcal U_{468}\longrightarrow\mathcal W_{243}\subset\mathbb F_3^6,
 \qquad
 \Psi(U_1,\ldots,U_6)=(-U_1,\ldots,-U_6)\pmod 3.
\]
Para \(w\in\mathbb F_3^6\), la aplicación
\[
 c(w)=(w,wA_W)\in C_W\subset\mathbb F_3^{12}
\]
es el codificador ternario fijado en la sección de Paley--Witt. Definimos
el lector de incidencia mínimo por
\[
 Q_{\min}=\operatorname{Supp}\circ c\circ\Psi:
 \mathcal U_{468}\longrightarrow\mathcal P([12]).
\]

Para una aplicación \(f:X\to Y\), denotemos por
\[
 \operatorname{Eq}(f)=\{(x,x')\in X^2:f(x)=f(x')\}
\]
su relación de equivalencia por fibras.

\begin{theorem}[Obstrucción del lector excepcional mínimo]
\label{p12-83-c17-s4:thm:obstruccion-lector-excepcional-minimo}
En el catálogo completo de 468 estados se verifica
\[
 \operatorname{Eq}(\Psi,Q_{\min})=\operatorname{Eq}(\Psi).
\]
En la notación abreviada habitual para particiones por fibras,
\[
 \ker(\Psi,Q_{\min})=\ker\Psi.
\]
Por consiguiente, \((\Psi,Q_{\min})\) no es inyectiva y el lector de
incidencia mínimo no distingue ningún par de estados que ya haya sido
identificado por \(\Psi\).
\end{theorem}

\begin{proof}
La factorización \(Q_{\min}=S\circ\Psi\), con
\(S=\operatorname{Supp}\circ c\), implica
\[
 \Psi(u)=\Psi(v)\Longrightarrow Q_{\min}(u)=Q_{\min}(v).
\]
La implicación recíproca entre las relaciones de equivalencia es inmediata,
pues \(\Psi\) es la primera componente de la aplicación diagonal. La
igualdad de relaciones se obtiene, por tanto, sin hipótesis estadísticas.
La enumeración exhaustiva confirma además que ambas particiones poseen el
mismo histograma:
\[
 132\,[1]+66\,[2]+18\,[3]+3\,[4]+6\,[5]+18\,[6],
\]
donde \(a\,[m]\) significa que existen \(a\) fibras de cardinal \(m\).
\end{proof}

La microfibra de \(w_\pi=010211\) proporciona el testigo distinguido. Sus
cinco regiones son
\[
\begin{split}
 &(501,614,498,169,272,272),\quad
 (810,923,870,169,272,272),\\
 &(810,983,810,169,272,272),\quad
 (870,923,810,169,272,272),\\
 &(870,923,810,418,674,521).
\end{split}
\]
Las cinco tienen la misma imagen visible y el mismo soporte excepcional
\[
 P_\pi=\{2,4,5,6,7,8,9,10,11\}.
\]
En el nivel de las 243 palabras aparecen 213 soportes distintos: 185 tienen
una preimagen, 27 tienen dos y uno tiene cuatro. La toma de soporte introduce,
por tanto, un segundo cociente además del impuesto por \(\Psi\).

El teorema delimita el lector mínimo, no la evaluación excepcional
enriquecida. Un lector capaz de separar las regiones de una misma palabra
debe actuar antes del cociente visible y utilizar al menos uno de los datos
que éste elimina: orientación, hoja, firma, retorno o estado de frontera. El
entrelazador dodecafásico desarrollado en la parte siguiente pertenece a esa
clase enriquecida y no contradice la obstrucción precedente.

La enumeración completa de las fibras produce el histograma anterior y
demuestra que la toma de soporte introduce un segundo cociente.


\section{Publicación arquimediana e incidencial}

Para cada nivel, sea \(\mathcal E_n\) un objeto incidencial y sean
\begin{equation}
 e_n:S_n\longrightarrow\mathcal E_n,
 \qquad
 \sigma_{n,m}:\mathcal E_n\longrightarrow\mathcal E_m
 \label{p12-83-c17-s5:eq:datos-lector-incidencial-u024}
\end{equation}
aplicaciones que satisfacen
\begin{equation}
 \sigma_{n,m}\circ e_n=e_m\circ r_{n,m}.
 \label{p12-83-c17-s5:eq:naturalidad-lector-incidencial-u024}
\end{equation}
La propiedad universal construye
\begin{equation}
 \operatorname{ev}_{\rm inc}:\Omega_\infty
 \longrightarrow\mathcal E_\infty,
 \qquad
 \mathcal E_\infty:=\varprojlim_n\mathcal E_n.
 \label{p12-83-c17-s5:eq:evaluacion-incidencial-limite-u024}
\end{equation}
La doble publicación es
\begin{equation}
 \mathcal P_\infty
 :=(\operatorname{ev}_{\mathbb R},\operatorname{ev}_{\rm inc}):
 \Omega_\infty\longrightarrow\mathbb R\times\mathcal E_\infty.
 \label{p12-83-c17-s5:eq:doble-publicacion-continuo-u024}
\end{equation}

\begin{figure}[H]
\centering
\begin{tikzcd}[column sep=10em,row sep=5em,
  cells={nodes={font=\large}}]
&\Omega_\infty
 \arrow[dl,"\operatorname{ev}_{\mathbb R}"']
 \arrow[dr,"\operatorname{ev}_{\rm inc}"]&\\
\mathbb R
 \arrow[dr,"\operatorname{coord}"']&&
\mathcal E_\infty
 \arrow[dl,"\operatorname{pub}_{\rm inc}"]\\
&\mathcal O&
\end{tikzcd}
\caption{Dos publicaciones de una misma historia. La conmutatividad en
\(\mathcal O\) no identifica los codominios \(\mathbb R\) y
\(\mathcal E_\infty\).}
\label{p12-83-c17-s5:fig:doble-publicacion-continuo-u024}
\end{figure}

La naturalidad demostrada en los
\cref{p12-83-c17-s2:thm:lector-excepcional-global-rev6,p12-83-c17-s2:thm:doble-realizacion-global-rev6}
se expresa, tras aplicar los lectores de publicación, mediante
\begin{equation}
 \operatorname{pub}_{\rm inc}\circ\operatorname{ev}_{\rm inc}
 =\operatorname{coord}\circ\operatorname{ev}_{\mathbb R},
 \label{p12-83-c17-s5:eq:compatibilidad-doble-publicacion-u024}
\end{equation}
y produce una comparación tipada de las dos publicaciones. La identidad
conserva el origen común sin construir una rama a partir de la otra: ambas
parten de \(\Omega_\infty\).


\section{Calibre de la publicación excepcional}

La segunda proyección no es un lector unario desprovisto de datos de
incidencia. Sea
\[
 \mathscr C_{\rm exc}
 =\mathscr C_{\rm ord}\times\mathscr C_{\rm Paley}
  \times\mathscr C_{\rm dod}\times\mathscr C_{\rm inc}
  \times\mathscr C_{\rm ret}
\]
el espacio de calibres excepcionales, cuyas coordenadas fijan,
respectivamente, el orden proyectivo, el calibre de Paley, la dodecada, el
alineamiento de incidencia y la cámara reticular. Se fija
\[
 \chi_{\rm exc}
 =(\chi_{\rm ord},\chi_{\rm Paley},\chi_{\rm dod},
   \chi_{\rm inc},\chi_{\rm ret})\in\mathscr C_{\rm exc}.
\]
El lector tipado tiene dominio
\[
 \mathfrak E:
 X_\infty^{\rm cal}\times\mathscr C_{\rm exc}
 \longrightarrow\mathbf{Exc},
 \qquad
 \mathfrak E_{\chi_{\rm exc}}(u):=\mathfrak E(u,\chi_{\rm exc}).
\]
En lo sucesivo, toda abreviatura \(\mathfrak E(u)\) significa
\(\mathfrak E_{\chi_{\rm exc}}(u)\) para este calibre explícitamente fijado.
La dependencia respecto de \(\chi_{\rm exc}\) no afecta a la existencia del
continuo ni a su publicación arquimediana: tipa la conservación excepcional
de la incidencia una vez construido el estado común.

\begin{falsifier}[Pérdida del calibre excepcional]
La comparación queda incompleta si se cambia el orden proyectivo, la dodecada,
el alineamiento de incidencia o la cámara reticular sin transportar las demás
coordenadas del calibre, o si se usa \(\mathfrak E\) como lector unario antes
de fijar \(\chi_{\rm exc}\). Tal fallo invalida esa comparación excepcional;
no modifica el cociente arquimediano ni regenera el continuo.
\end{falsifier}


La segunda lectura no nace del número real. Sea
\[
 \mathfrak E:
 \mathscr O_{\mathbb R,\leftrightarrow}^{\rm HMT}
 \longrightarrow\mathbf{Exc}
\]
el lector continuo de incidencia, con codominio de Hausdorff, que prolonga
la cadena de Paley--Witt, Golay,
Leech, \(V^\natural\), Monstruo y Moonshine. Ambas aplicaciones parten de
la misma historia:
\begin{equation}
 \mathfrak D(u)=(\mathfrak R(u),\mathfrak E(u)).
 \label{p12-83-c17-s7:hmtch:eq-doble-proyeccion-two-way}
\end{equation}
La primera contrae las fibras compatibles hasta su media arquimediana; la
segunda conserva incidencia, orientación y memoria de borde.

Fijado el origen del calendario dodecafásico, sea
\[
 R_+:=\{k\in\mathbb Z:1+(k\bmod 9)\in\{1,4,7\}\}
\]
el conjunto biinfinito de retornos completos \((\tau=+1)\). Es sindético:
sus huecos son uniformemente acotados. Para
\(u=(m,(x_k)_{k\in\mathbb Z})\) definimos
\[
 \operatorname{Ret}_+(u)
 :=\bigl(m,(x_k)_{k\in R_+}\bigr)
 \in\mathbb Z\times\prod_{k\in R_+}X,
\]
pero cada \(x_k\) es aquí el estado entero
\eqref{p12-83-c19-s1:hmtch:eq-celula-trit-completa} con hoja, carry, ruta e incidencia,
no el signo aislado.

\begin{proposition}[Muestreo sin pérdida del registro genealógico de retorno]
\label{p12-83-c17-s7:hmtch:prop-retorno-reconstruye}
La aplicación \(\operatorname{Ret}_+\), con la topología de subespacio en el
producto, es un homeomorfismo sobre su imagen.
En consecuencia existe una única aplicación continua
\(\mathfrak E_+\) sobre esa imagen tal que
\begin{equation}
 \mathfrak E=\mathfrak E_+\circ\operatorname{Ret}_+.
 \label{p12-83-c17-s7:hmtch:eq-factorizacion-excepcional-retorno}
\end{equation}
\end{proposition}

\begin{proof}
Entre dos retornos consecutivos hay un número uniformemente acotado de pasos
del calendario. El estado completo en un retorno contiene la transición y
el registro genealógico que determinan los pasos intermedios; hacia la otra dirección se
usa la flecha inversa de la completación grupoidal. Así se reconstruye toda
la órbita. La aplicación es continua y, al restringirla a una carta
compacta, una inyección continua en un espacio de Hausdorff es un
homeomorfismo sobre su imagen. La coordenada entera se conserva
explícitamente. La factorización se define transportando \(\mathfrak E\) por
la inversa y es continua carta a carta.
\end{proof}

Por eso el \(+1\) participa en el corredor excepcional como memoria completa
de retorno: el muestreo de estados enteros no pierde la salida excepcional.
Esta proposición demuestra recuperación y factorización, no una generación
del corredor a partir del escalar \(+1\) aislado. La salida arquimediana y la
salida excepcional son dos resultantes coordinadas del mismo antecedente.
